\section{矩阵与狄拉克符号}
首先，考虑算符 $X$。我们需要构造一个矩阵来等效描述它，众所周知，矩阵的元素必须是标量（数）而非矢量，因此我们可以将单位算符作用于 $X$：
\begin{equation}
	X = \sum_{a'} \sum_{a''} \ket{a'} \bra{a'} X \ket{a''} \bra{a''}
\end{equation}

接下来我们构造标量 $\bra{a'} X \ket{a''}$，只要能为态矢量构造出与 $\bra{a'} X \ket{a''}$ 匹配的恰当矩阵形式，就可以将该标量作为矩阵的元素，形式如下：
\begin{equation}
	X \doteq \begin{pmatrix}
		\bra{a_{(1)}} X \ket{a_{(1)}} & \bra{a_{(1)}} X \ket{a_{(2)}} & \cdots \\
		\bra{a_{(2)}} X \ket{a_{(1)}} & \bra{a_{(2)}} X \ket{a_{(2)}} & \cdots \\
		\vdots & \vdots & \ddots
	\end{pmatrix}
\end{equation}

为了让态矢量与算符 $X$ 的形式匹配，我们回顾算符 $X$ 的定义——态矢量空间中的基本方程：
\begin{equation}
	\ket{\beta} = X \ket{\alpha}
\end{equation}

将单位算符作用于该式两侧：
\begin{equation}
	\bra{a_l} \ket{\beta} = \sum_m \bra{a_l} X \ket{a_m} \bra{a_m} \ket{\alpha}
\end{equation}

若我们将 $\bra{a_l} \ket{\beta}$ 和 $\bra{a_m} \ket{\alpha}$ 排列为列矩阵，并定义 $\ket{\beta}$ 和 $\ket{\alpha}$ 的新表示形式：
\begin{equation}
	|\beta\rangle \doteq \begin{pmatrix}
		\langle a_1 | \beta \rangle \\
		\langle a_2 | \beta \rangle \\
		\vdots
	\end{pmatrix}, \quad
	|\alpha\rangle \doteq \begin{pmatrix}
		\langle a_1 | \alpha \rangle \\
		\langle a_2 | \alpha \rangle \\
		\vdots
	\end{pmatrix}
\end{equation}

此时我们会惊奇地发现，上式恰好对应矩阵乘法运算。因此，若我们按上述方式定义 $X$、$|\alpha\rangle$ 和 $|\beta\rangle$ 的矩阵表示，它们将完全契合狄拉克态矢量理论。

类似地，我们可以推导出左矢的恰当表示形式：
\begin{equation}
	\langle \beta| = \langle \alpha| X
\end{equation}
\begin{equation}
	\langle \beta| a'\rangle = \sum_m \langle \alpha| a_m \rangle \langle a_m | X | a' \rangle
\end{equation}
\begin{equation}
	\langle \beta| \doteq \left( \langle a_1 | \beta \rangle \ \langle a_2 | \beta \rangle \ \cdots \right), \quad
	\langle \alpha| \doteq \left( \langle a_1 | \alpha \rangle \ \langle a_2 | \alpha \rangle \ \cdots \right)
\end{equation}

由此我们可以大胆得出结论：若将态矢量和算符替换为对应的矩阵形式，同时保持标量不变，那么态矢量理论中的所有方程形式都将保持不变。接下来我们可以通过内积验证该理论的正确性：
\begin{equation}
	\langle \beta | \alpha \rangle = \sum_i \langle \beta | a_i \rangle \langle a_i | \alpha \rangle
\end{equation}

在矩阵表示下：
\begin{equation}
	\langle \beta | \alpha \rangle \doteq \left( \langle a_1 | \beta \rangle \ \langle a_2 | \beta \rangle \ \cdots \right) \begin{pmatrix}
		\langle a_1 | \alpha \rangle \\
		\langle a_2 | \alpha \rangle \\
		\vdots
	\end{pmatrix} = \sum_a \langle \beta | a \rangle \langle a | \alpha \rangle
\end{equation}

我们发现，无论采用哪种表示形式，$\langle \beta | \alpha \rangle$ 的结果（标量）始终一致，这证明了矩阵表示的正确性。

\subsection{基变换}
\subsubsection{变换关系}

在量子力学中，我们可以选择不同的基右矢，以 $|a\rangle$ 和 $|b\rangle$ 为例。因此，要研究基变换问题，首先需要分析变换算符 $U$：
\begin{equation}
	|b_i \rangle = U |a_i \rangle
\end{equation}
若这些基右矢满足正交归一性，则可推得：
\begin{equation}
	U = \sum_{l} \ket{b_{l}} \bra{a_{l}}
\end{equation}
\begin{equation}
	UU^{\dagger} = 1
\end{equation}

我们同样可以写出 $U$ 的矩阵表示：
\begin{equation}
	\bra{a_{l}} U \ket{a_{m}} = \braket{a_{l}}{b_{m}}
\end{equation}

接下来分析算符 $X$ 在新基下的变换形式：

首先，在矩阵表示下：
\begin{align}
	\bra{b_{l}} X \ket{b_{m}} &= \sum_{p} \sum_{q} \braket{b_{l}}{a_{p}} \bra{a_{p}} X \ket{a_{q}} \braket{a_{q}}{b_{m}} \\
	&= \sum_{p} \sum_{q} \bra{a_{l}} U^{\dagger} \ket{a_{p}} \bra{a_{p}} X \ket{a_{q}} \bra{a_{q}} U \ket{a_{m}}
\end{align}

这一形式恰好对应矩阵乘法，因此我们得到算符 $X$ 的变换关系：
\begin{equation}
	X' = U^{\dagger} X U
\end{equation}

\subsubsection{对角化}

借助矩阵表示，我们可以利用线性代数中的方法求解量子力学中的本征方程。考虑算符 $B$ 的本征方程：
\begin{equation}
	B \ket{b_{l}} = b_{l} \ket{b_{l}}
\end{equation}

已知 $B$ 在基 $\ket{a}$ 下的矩阵表示，将其代入本征方程可得：
\begin{equation}
	\sum_{k} \bra{a_{m}} B \ket{a_{k}} \braket{a_{k}}{b_{l}} = b_{l} \braket{a_{m}}{b_{l}}
\end{equation}

若定义 $C_{m l} = \braket{a_{m}}{b_{l}}$，则上式可写为：
\begin{equation}
	\begin{pmatrix}
		B_{11} & B_{12} & \cdots \\
		B_{21} & B_{22} & \cdots \\
		\vdots & \vdots & \ddots
	\end{pmatrix}
	\begin{pmatrix}
		C_{1 l} \\
		C_{2 l} \\
		\vdots
	\end{pmatrix}
	= b_{l}
	\begin{pmatrix}
		C_{1 l} \\
		C_{2 l} \\
		\vdots
	\end{pmatrix}
\end{equation}

由此，我们可以利用线性代数中的方法（求解 $\det(B - \lambda I) = 0$）得到本征值 $b_{l}$ 和列向量 $\begin{pmatrix} C_{1 l} \\ C_{2 l} \\ \vdots \end{pmatrix}$。注意到 $C_{m l} = \braket{a_{m}}{b_{l}}$ 意味着该列向量正是本征右矢 $\ket{b_{l}}$ 在基 $\ket{a}$ 下的矩阵表示。至此，我们通过矩阵表示得到了本征值 $b_{l}$ 和本征右矢 $\ket{b_{l}}$.

接下来分析如何将基 $\ket{a}$ 下的矩阵 $B$ 对角化：显然，在 $B$ 的本征基 $\ket{b}$ 下，$B$ 是对角矩阵：
\begin{equation}
	\langle b_m | B | b_l \rangle = b_l \delta_{ml}
\end{equation}

因此，我们只需通过基变换即可实现对角化：
\begin{equation}
	B_b = U^\dagger B_a U \quad B_a = U B_b U^{-1}
\end{equation}

那么矩阵 $U$ 的具体形式是什么？幸运的是，$C_{m l} = \langle a_m | b_l \rangle$ 恰好是矩阵 $U$ 第 $m$ 行第 $l$ 列的元素——将所有列向量 $\begin{pmatrix} C_{1 1} \\ C_{2 1} \\ \vdots \end{pmatrix}$ 按列排列为方阵 $\begin{pmatrix} C_{1 1} & C_{1 2} & \cdots \\ C_{2 1} & C_{2 2} & \cdots \\ \vdots & \vdots & \ddots \end{pmatrix}$，该方阵即为 $U$。

简言之，矩阵对角化的步骤为：首先求解本征方程得到本征值和本征右矢；然后将所有本征右矢的矩阵表示按列排列为方阵 

\subsection{基本对易关系}
坐标与动量的正则对易关系
\begin{equation}
	[x_i, p_j] = i\hbar \delta_{ij}
\end{equation}

角动量分量的对易关系
\begin{equation}
	[L_i, L_j] = i\hbar \epsilon_{ijk} L_k
\end{equation}

角动量平方与分量的对易关系
\begin{equation}
	[L^2, L_i] = 0
\end{equation}

角动量与位置算符
\begin{equation}
	[L_i, x_j] = i\hbar \epsilon_{ijk} x_k
\end{equation}

角动量与动量算符
\begin{equation}
	[L_i, p_j] = i\hbar \epsilon_{ijk} p_k
\end{equation}

Pauli 矩阵的对易关系
\begin{equation}
	[\sigma_i, \sigma_j] = 2i \epsilon_{ijk} \sigma_k
\end{equation}

对任意 \(i,j \in \{1,2,3\}\), Pauli 矩阵的乘积可写为
\[
\sigma_i \sigma_j = \delta_{ij} I + i\sum_{k=1}^3 \epsilon_{ijk}\sigma_k
\]

自旋角动量分量的对易关系
\begin{equation}
	[S_i, S_j] = i\hbar \epsilon_{ijk} S_k
\end{equation}

对易关系的基本性质
\begin{equation}
	[A, B + C] = [A, B] + [A, C], \quad [A + B, C] = [A, C] + [B, C] 
\end{equation}
\begin{equation}
	[A, BC] = [A, B]C + B[A, C], \quad [AB, C] = A[B, C] + [A, C]B
\end{equation}
\begin{equation}
	[A, [B, C]] + [B, [C, A]] + [C, [A, B]] = 0
\end{equation}